%% LyX 2.3.6.2 created this file.  For more info, see http://www.lyx.org/.
%% Do not edit unless you really know what you are doing.
\documentclass[english]{article}
\usepackage{mathptmx}
\renewcommand{\familydefault}{\rmdefault}
\usepackage[T1]{fontenc}
\usepackage[latin9]{inputenc}
\usepackage{geometry}
\geometry{verbose,tmargin=1in,bmargin=1in,lmargin=1in,rmargin=1in}
\usepackage{babel}
\usepackage{esint}
\usepackage[unicode=true,pdfusetitle,
 bookmarks=true,bookmarksnumbered=false,bookmarksopen=false,
 breaklinks=false,pdfborder={0 0 1},backref=false,colorlinks=false]
 {hyperref}

\makeatletter
%%%%%%%%%%%%%%%%%%%%%%%%%%%%%% User specified LaTeX commands.
\newcommand{\vect}[1]{\mbox{\boldmath $#1$}}

\makeatother

\begin{document}
\title{Alpha-particle loss fraction metric}
\author{Matt Landreman}

\maketitle
The fraction of lost alpha particles is
\begin{equation}
f=\frac{N}{D},\label{eq:f_def}
\end{equation}
where
\begin{equation}
N=\int d^{3}x\int d^{3}v\,S\Theta,
\end{equation}
\begin{equation}
D=\int d^{3}x\int d^{3}v\,S,
\end{equation}
$\int d^{3}x$ is an integral over the plasma volume, $S$ is the
relative local rate at which alpha particles are born, and $\Theta$
is an indicator function which is 1 if the alpha born at this location
is lost before thermalization, 0 otherwise. We will compute $\Theta$
by following contours of the adiabatic invariant $J=\int v_{||}d\ell$.
We will assume particles conserve $J$ aside from transitions between
trapping classes, at which point $J$ changes discontinuously. Particles
can be followed through a transition by considering all the possible
classes they may transition into. If a point in phase space can reach
the plasma boundary through transitions in this way and otherwise
following constant-$J$ contours, we will set $\Theta=1$, otherwise
$\Theta=0$. For a particle that never bounces, we will consdier $\Theta=0$.

We only need to know $S$ up to a constant because if it is multiplied
by any scale factor, the scale factor cancels out between numerator
and denominator.

The source has the form
\begin{equation}
S=\delta\left(W-W_{0}\right)h\left(\rho\right)
\end{equation}
where
\begin{equation}
W=\frac{v^{2}}{2}
\end{equation}
is the specific kinetic energy, $W_{0}=3.5$ MeV$/m_{\alpha}$ is
the birth specific energy, $\rho$ is a surface label that ranges
over $[0,1]$, and $h$ is some function of it. 

We convert the spatial variables to coordinates $\left(\rho,\alpha,\zeta\right)$
where
\begin{equation}
\alpha=\theta-\iota\zeta
\end{equation}
is a field line label and $\left(\theta,\zeta\right)$ are any straight-field-line
coordinates. We have
\begin{equation}
\mathbf{B}=\nabla\psi\times\nabla\theta+\iota\nabla\zeta\times\nabla\psi=\nabla\psi\times\nabla\alpha,
\end{equation}
where $2\pi\psi$ is the toroidal flux. The Jacobian of the coordinate
transformation is
\begin{equation}
\sqrt{g}=\frac{1}{\nabla\rho\times\nabla\alpha\cdot\nabla\zeta}=\frac{d\psi}{d\rho}\frac{1}{\nabla\psi\times\nabla\alpha\cdot\nabla\zeta}=\frac{d\psi}{d\rho}\frac{1}{\mathbf{B}\cdot\nabla\zeta}
\end{equation}
so
\begin{equation}
\int d^{3}x=\int_{0}^{1}d\rho\int_{0}^{2\pi}d\alpha\int_{0}^{2\pi}d\zeta\left|\frac{d\psi}{d\rho}\frac{1}{\mathbf{B}\cdot\nabla\zeta}\right|.
\end{equation}

For the denominator $D$, the integrand is isotropic so we can work
in terms of spherical coordinates in velocity space:
\begin{equation}
\int d^{3}v=4\pi\int_{0}^{\infty}dv\;v^{2}=4\pi\int_{0}^{\infty}dW\,\sqrt{2W}.
\end{equation}
Thus, the denominator is
\begin{equation}
D=4\pi\sqrt{2W_{0}}\int_{0}^{1}d\rho\,h\int_{0}^{2\pi}d\alpha\int_{0}^{2\pi}d\zeta\left|\frac{d\psi}{d\rho}\frac{1}{\mathbf{B}\cdot\nabla\zeta}\right|.
\end{equation}
If we notice
\begin{equation}
\frac{dV}{d\rho}=\frac{d}{d\rho}\left(\int d^{3}x\right)=\int_{0}^{2\pi}d\alpha\int_{0}^{2\pi}d\zeta\left|\frac{d\psi}{d\rho}\frac{1}{\mathbf{B}\cdot\nabla\zeta}\right|
\end{equation}
then another expression for $D$ is
\begin{equation}
D=4\pi\sqrt{2W_{0}}\int_{0}^{1}d\rho\,h\frac{dV}{d\rho}.
\end{equation}

Now we turn to the numerator $N$. As shown in the appendix, the velocity
integral can be written
\begin{equation}
\int d^{3}v=\sum_{\sigma}\int_{0}^{2\pi}d\vartheta\int_{0}^{\infty}dW\int_{0}^{W/B}d\mu\frac{B}{\left|v_{||}\right|},\label{eq:d3v}
\end{equation}
where $\sigma=\mbox{sign}\left(v_{||}\right)=\pm1$, $\vartheta$
is the gyrophase, $W=v^{2}/2$ is again the specific kinetic energy,
and $\mu=v_{\perp}^{2}/\left(2B\right)$ is the magnetic moment. Since
we are using guiding-center calculations to determine $\Theta$, nothing
in the integrand depends on $\vartheta$. Therefore
\begin{equation}
N=2\pi\int_{0}^{1}d\rho\,h\left|\frac{d\psi}{d\rho}\right|\int_{0}^{2\pi}d\alpha\int_{0}^{2\pi}d\zeta\sum_{\sigma}\int_{0}^{W_{0}/B}d\mu\frac{\Theta}{\left|v_{||}\mathbf{b}\cdot\nabla\zeta\right|}.
\end{equation}
We can pull the $\mu$ integral outside the others as long as the
ranges of integration are adjusted to cover the same region of phase
space:
\begin{equation}
N=2\pi\sum_{\sigma}\sum_{\mathrm{wells}}\int_{0}^{W_{0}/B_{\min}}d\mu\int_{0}^{1}d\rho\,h\left|\frac{d\psi}{d\rho}\right|\int d\alpha\int d\zeta\frac{\Theta}{\left|v_{||}\mathbf{b}\cdot\nabla\zeta\right|}.
\end{equation}
Here, the notation for the range of the $\alpha$ and $\zeta$ integrals
has been left vague deliberately - these integrals and the sum over
wells should be done such that phase space is weighted and counted
properly, without over-counting or under-counting any wells.

Now let us change variables from $\mu$ to $\lambda_{n}$, where $\lambda_{n}$
is defined via
\begin{equation}
B_{\mathrm{bounce}}=B_{\min}+\left(B_{\max}-B_{\min}\right)\lambda_{n}.
\end{equation}
To make this change of variables we first note
\begin{equation}
\frac{v^{2}}{2}=\frac{v_{||}^{2}}{2}+\frac{v_{\perp}^{2}}{2}
\end{equation}
so
\begin{equation}
W=\frac{v_{||}^{2}}{2}+\mu B.
\end{equation}
At a bounce point, we have $v_{||}\to0$ and $B\to B_{\mathrm{bounce}}$,
and we always have $W=W_{0}$, so
\begin{equation}
W_{0}=\mu B_{\mathrm{bounce}}.
\end{equation}
Hence,
\begin{equation}
W_{0}=\left[B_{\min}+\left(B_{\max}-B_{\min}\right)\lambda_{n}\right]\mu.
\end{equation}
Differentiating,
\begin{equation}
0=\left[B_{\min}+\left(B_{\max}-B_{\min}\right)\lambda_{n}\right]d\mu+\mu\left(B_{\max}-B_{\min}\right)d\lambda_{n},
\end{equation}
so
\begin{equation}
\frac{d\mu}{d\lambda_{n}}=-\mu\frac{\left(B_{\max}-B_{\min}\right)}{B_{\min}+\left(B_{\max}-B_{\min}\right)\lambda_{n}}=-W_{0}\frac{\left(B_{\max}-B_{\min}\right)}{\left[B_{\min}+\left(B_{\max}-B_{\min}\right)\lambda_{n}\right]^{2}}.
\end{equation}
We only need to integrate over $\lambda_{n}\in\left[0,1\right]$ because
passing particles are not lost and so have $\Theta=0$. Thus,
\begin{equation}
N=2\pi W_{0}\sum_{\sigma}\sum_{\mathrm{wells}}\int_{0}^{1}d\lambda_{n}\frac{\left(B_{\max}-B_{\min}\right)}{\left[B_{\min}+\left(B_{\max}-B_{\min}\right)\lambda_{n}\right]^{2}}\int_{0}^{1}d\rho\,h\left|\frac{d\psi}{d\rho}\right|\int d\alpha\int d\zeta\frac{\Theta}{\left|v_{||}\mathbf{b}\cdot\nabla\zeta\right|}.
\end{equation}
Defining the birth speed $v_{0}$ using 
\begin{equation}
W_{0}=\frac{v_{0}^{2}}{2}
\end{equation}
and noting
\begin{equation}
\left|v_{||}\right|=v_{0}\sqrt{1-B/B_{\mathrm{bounce}}}
\end{equation}
then
\begin{equation}
N=\pi v_{0}\sum_{\sigma}\sum_{\mathrm{wells}}\int_{0}^{1}d\lambda_{n}\frac{\left(B_{\max}-B_{\min}\right)}{B_{\mathrm{bounce}}^{2}}\int_{0}^{1}d\rho\,h\left|\frac{d\psi}{d\rho}\right|\int d\alpha\int d\zeta\frac{\Theta}{\left|\mathbf{b}\cdot\nabla\zeta\right|\sqrt{1-B/B_{\mathrm{bounce}}}}.
\end{equation}

An expression for the alpha loss fraction is then
\begin{equation}
f=\frac{\pi v_{0}\sum_{\sigma}\sum_{\mathrm{wells}}\int_{0}^{1}d\lambda_{n}\frac{\left(B_{\max}-B_{\min}\right)}{B_{\mathrm{bounce}}^{2}}\int_{0}^{1}d\rho\,h\left|\frac{d\psi}{d\rho}\right|\int d\alpha\int d\zeta\frac{\Theta}{\left|\mathbf{b}\cdot\nabla\zeta\right|\sqrt{1-B/B_{\mathrm{bounce}}}}}{4\pi\sqrt{2W_{0}}\int_{0}^{1}d\rho\,h\int_{0}^{2\pi}d\alpha\int_{0}^{2\pi}d\zeta\left|\frac{d\psi}{d\rho}\frac{1}{\mathbf{B}\cdot\nabla\zeta}\right|}.
\end{equation}
If $\Theta$ is evaluated using $J$ contours, as we are proposing,
then $\Theta$ will be independent of $\sigma$ so we can replace
$\sum_{\sigma}\to2$. Then, a final expression for the loss fraction
becomes
\begin{equation}
f=\frac{\sum_{\mathrm{wells}}\int_{0}^{1}d\lambda_{n}\frac{\left(B_{\max}-B_{\min}\right)}{B_{\mathrm{bounce}}^{2}}\int_{0}^{1}d\rho\,h\left|\frac{d\psi}{d\rho}\right|\int d\alpha\int d\zeta\frac{\Theta}{\left|\mathbf{b}\cdot\nabla\zeta\right|\sqrt{1-B/B_{\mathrm{bounce}}}}}{2\int_{0}^{1}d\rho\,h\int_{0}^{2\pi}d\alpha\int_{0}^{2\pi}d\zeta\left|\frac{d\psi}{d\rho}\frac{1}{\mathbf{B}\cdot\nabla\zeta}\right|}.
\end{equation}

As already emphasized, the sum over wells and the integration ranges
should be interpreted as covering phase space once, and not over-counting
or under-counting any wells, such that the overall weight on phase
space is equivalent to the original weight in equation (\ref{eq:f_def}).

Nothing in the calculation changes if we use $\theta$ as an integration
coordinate instead of $\alpha$ - we would just need to evaluate the
integrands as functions of $\theta$ then instead of as functions
of $\alpha$. We could also evaluate the numerator with an integral
$\int d\alpha$ while evaluating the denominator with an integral
$\int d\theta$.

\section{Appendix: velocity-space integral}

Here we derive (\ref{eq:d3v}). First observe that $\left(v_{\perp},\vartheta,v_{||}\right)$
are cylindrical coordinates for velocity space, so
\begin{equation}
\int d^{3}v=\int_{0}^{2\pi}d\vartheta\int_{-\infty}^{\infty}dv_{||}\int_{0}^{\infty}dv_{\perp}v_{\perp}.
\end{equation}
Then using $\mu=v_{\perp}^{2}/\left(2B\right)$, 
\begin{equation}
2B\;d\mu=2v_{\perp}\;dv_{\perp}
\end{equation}
so
\begin{equation}
\int d^{3}v=\int_{0}^{2\pi}d\vartheta\int_{-\infty}^{\infty}dv_{||}\int_{0}^{\infty}d\mu\,B.
\end{equation}
We now write the energy as
\begin{equation}
W=\frac{v_{||}^{2}}{2}+\mu B
\end{equation}
so at fixed $\mu$ and position
\begin{equation}
dW=v_{||}dv_{||}.
\end{equation}
Hence
\begin{eqnarray}
\int d^{3}v & = & \sum_{\sigma}\int_{0}^{2\pi}d\vartheta\int_{0}^{\infty}d\mu\int_{\mu B}^{\infty}dW\frac{B}{\left|v_{||}\right|}\label{eq:d3v_integral}\\
 & = & \sum_{\sigma}\int_{0}^{2\pi}d\vartheta\int_{0}^{\infty}dW\int_{0}^{W/B}d\mu\frac{B}{\left|v_{||}\right|}.\nonumber 
\end{eqnarray}

\end{document}
